\chapter*{Abstract}
\addcontentsline{toc}{chapter}{Abstract}

The evidence for dark matter comes from many independent probes at different scales. Galactic rotation curves, the motions of galaxies within clusters, gravitational lensing, and the anisotropies of the cosmic microwave background all require a form of matter that outweighs ordinary baryons by roughly a factor of five and interacts with light only through gravity. These probes, however, are all gravitational in nature, and they say little about what dark matter actually is. If it is made of weakly interacting massive particles, it could annihilate or decay into Standard Model states, and gamma rays are especially attractive messengers: undeflected by magnetic fields, they carry the spectral signature of the interaction that produced them. For more than fifteen years the Large Area Telescope on board the Fermi Gamma-ray Space Telescope (Fermi-LAT) has surveyed the GeV sky where these signatures are expected, and its measurements are at the core of this thesis.

In this thesis, we argue that progress in indirect dark matter searches now depends on statistical and machine-learning methods capable of recovering signals in noise-dominated environments. The work follows an arc from the Galactic Center outward to the cosmic web, tracking the expected signal into regimes of progressively lower signal-to-noise ratio. To this end, we draw on traditional statistical methods, generative modeling, implicit likelihood methods, and deep learning, applied to Fermi-LAT data and to forecasts for future instruments.

The Galactic Center hosts the Galactic Center Excess (GCE), a surplus of GeV photons that may signal dark matter annihilation or an unresolved population of millisecond pulsars (MSPs). To test the pulsar hypothesis we measured the gamma-ray luminosity function of MSPs in Milky Way globular clusters, finding a mean luminosity $\langle L_\gamma\rangle \sim (1\text{--}8)\times10^{33}$ erg/s between 0.1 and 100 GeV and a broad log-normal spread, $\sigma_L \sim 1.4\text{--}2.8$. Had the GCE been produced by pulsars with these properties, Fermi-LAT would already have resolved roughly 17 to 37 of them individually, whereas only three candidates are known --- a tension that reignites the debate over whether the excess comes from pulsars or from dark matter. 

We then turned to fainter targets: dark matter subhalos potentially hidden among the Fermi-LAT sources that lack an astrophysical counterpart. We built what is, to the best of our knowledge, the first generative model of these unassociated sources at latitudes above ten degrees: a statistical mixture of known astrophysical classes and a possible subhalo component, trained directly on this population and corrected for potential prior and covariate shifts. 
We found no significant subhalo contribution and set 95\% confidence upper limits on the annihilation cross section in the $b\bar{b}$ channel for dark matter masses between 10 GeV and 1 TeV.

Below the Fermi-LAT detection threshold, individual sources blur together, and the object of study becomes the gamma-ray source population as a whole --- astrophysical and potential dark matter sources alike. We trained a convolutional neural network on about $10^6$ synthetic Fermi-LAT sky maps to reconstruct the source-count distribution $dN/dS$ of extragalactic sources from fourteen years of data between 1 and 10 GeV. The reconstruction reaches fluxes a factor of about 50 below the Fermi-LAT threshold, agrees with existing catalogs in the resolved regime, and extends as $dN/dS \propto S^{-2}$ down to roughly $5\times10^{-12}$ cm$^{-2}$ s$^{-1}$. 
Building on this distribution, we developed a probabilistic cataloging framework that compares the observed sky with simulations and assigns each direction a likelihood of hosting a sub-threshold source. This framework identifies roughly 50\% more candidate source directions than the 4FGL-DR3 catalog, and delivers a public catalog particularly suited to cross-correlation studies.

At the largest scales, the annihilation signal follows the matter distribution itself, so cross-correlating gamma-ray anisotropies with galaxy positions can separate a cosmological dark matter component from isotropic backgrounds. We forecast the sensitivity of the Cherenkov Telescope Array Observatory (CTAO) to this signal. With a dense low-redshift galaxy catalog such as 2MASS and about 50 hours of observation, the cross-correlation reaches sensitivities to both annihilating and decaying dark matter comparable to those of dwarf-galaxy and galaxy-cluster analyses.

Taken together, these results trace a coherent path through the gamma-ray sky, from the brightest anticipated signal to the faint imprint of the cosmic web. None of these searches returns a confirmed detection, and the particle identity of dark matter remains an open problem. The next steps in indirect detection, we argue, will require better instruments, more data, and sharper statistical tools to read them, all advancing together. This thesis contributes to the third of these fronts, working toward new statistical tools for the field: population models that make sub-threshold sources statistically accessible and inference frameworks robust to the mismatch between simulations and real data, together with probabilistic catalogs and cross-correlation analyses. The hope is that this groundwork will help extend these searches as current gamma-ray surveys continue to grow and next-generation observatories come online.
